% !TeX program = lualatex
\documentclass[12pt,aspectratio=169]{beamer}
\usepackage[utf8]{inputenc}
\usepackage[T1]{fontenc}
\usepackage{amsmath,amsfonts,amssymb}
\usepackage{graphicx}
\usepackage{tikz}
\usepackage{xcolor}
\usepackage{hyperref}
\usepackage{anyfontsize}
\usepackage{lmodern}

% ====== MAKE ALL FONTS SMALLER ======
\renewcommand{\tiny}{\fontsize{5}{6}\selectfont}
\renewcommand{\scriptsize}{\fontsize{6}{7}\selectfont}
\renewcommand{\footnotesize}{\fontsize{7}{8}\selectfont}
\renewcommand{\small}{\fontsize{8}{9}\selectfont}
\renewcommand{\normalsize}{\fontsize{9}{10}\selectfont}
\renewcommand{\large}{\fontsize{10}{11}\selectfont}
\renewcommand{\Large}{\fontsize{11}{13}\selectfont}
\renewcommand{\LARGE}{\fontsize{13}{15}\selectfont}
\renewcommand{\huge}{\fontsize{15}{17}\selectfont}
\renewcommand{\Huge}{\fontsize{18}{21}\selectfont}

% ====== COLOR THEME ======
\definecolor{redcorrect}{RGB}{200,30,30}
\definecolor{bluecorrect}{RGB}{30,80,200}
\definecolor{greencheck}{RGB}{0,150,50}
\definecolor{lightgray}{RGB}{245,245,245}
\definecolor{darkgray}{RGB}{50,50,50}
\definecolor{softyellow}{RGB}{255,248,200}
\definecolor{infoblue}{RGB}{230,240,252}
\definecolor{steppurple}{RGB}{150,50,200}
\definecolor{deepblue}{RGB}{20,60,150}
\definecolor{orangewarning}{RGB}{255,140,0}

\setbeamercolor{title}{fg=white,bg=redcorrect}
\setbeamercolor{frametitle}{fg=white,bg=redcorrect}
\setbeamercolor{block title}{fg=white,bg=bluecorrect}
\setbeamercolor{block body}{fg=darkgray,bg=lightgray}
\setbeamercolor{normal text}{fg=darkgray}
\usecolortheme{default}

% ====== CUSTOM COMMANDS ======
\newcommand{\correct}[1]{\textcolor{greencheck}{\textbf{\checkmark}} \textcolor{greencheck}{#1}}
\newcommand{\highlightred}[1]{\textcolor{redcorrect}{\textbf{#1}}}
\newcommand{\highlightblue}[1]{\textcolor{bluecorrect}{\textbf{#1}}}
\newcommand{\highlightorange}[1]{\textcolor{orangewarning}{\textbf{#1}}}
\newcommand{\tipbox}[1]{\fcolorbox{bluecorrect}{softyellow}{\parbox{0.88\textwidth}{\centering #1}}}
\newcommand{\steptitle}[1]{\textcolor{steppurple}{\textbf{#1}}}
\newcommand{\keypoint}[1]{\textcolor{deepblue}{\textbf{#1}}}
\newcommand{\redtitle}[1]{\textcolor{redcorrect}{\textbf{#1}}}

\newcommand{\stepbox}[1]{%
	\begin{center}
		\fcolorbox{darkgray}{white}{%
			\parbox{0.90\textwidth}{\vspace{0.1cm} #1 \vspace{0.1cm}}%
		}
	\end{center}
}

\begin{document}

% ================================================================
% SLIDE 1: TITLE
% ================================================================
\begin{frame}
\centering
\vspace{0.8cm}
{\fontsize{20}{24}\selectfont \textbf{Money Laundering and Financial Crime}}\\[0.2cm]
{\fontsize{15}{18}\selectfont Complete Tutorial Guide — All Topics}\\[0.4cm]
{\fontsize{12}{14}\selectfont Step-by-Step Learning for Every Subject}\\[0.4cm]
{\fontsize{10}{12}\selectfont Compliance Training Department}
\vspace{0.3cm}
\rule{5cm}{1pt}
\end{frame}

% ================================================================
% SLIDE 2: TUTORIAL 1 — OPEN-SOURCE RESEARCH
% ================================================================
\begin{frame}
\frametitle{\fontsize{14}{17}\selectfont TUTORIAL 1: Open-Source Research in Investigations}

\small
\begin{block}{\fontsize{11}{13}\selectfont Topic: Transaction Monitoring and Investigation}
\footnotesize
When reviewing alerts, analysts often need to understand unfamiliar entities. Open-source research is a critical tool for this.
\end{block}

\vspace{0.03cm}

\stepbox{\footnotesize
\steptitle{Step 1:} Understand the scenario. An analyst notices wire transfers to an unfamiliar entity in a sanctioned jurisdiction. The customer profile does not list dealings in that region. \\
\steptitle{Step 2:} Recognize the need for information. The analyst needs to understand who the entity is and why the customer is transacting with them. \\
\steptitle{Step 3:} Understand the correct source. \textbf{Open-source research into the entity} is the most helpful information source. \\
\steptitle{Step 4:} See why this is the best approach. Open-source research can reveal: \\
\quad - The nature of the entity (company, individual, shell company) \\
\quad - Connections to sanctioned individuals or entities \\
\quad - Public information about the entity's activities \\
\steptitle{Step 5:} Understand what is NOT the best answer. \textbf{Regulatory advisories about common evasion techniques} are general information, not specific to this entity. The specific entity research is more valuable.}

\vspace{0.08cm}
\tipbox{\footnotesize \textbf{CORE LESSON:} When investigating unfamiliar entities, open-source research provides specific information about the entity itself. This is more valuable than general regulatory advisories.}
\end{frame}

% ================================================================
% SLIDE 3: TUTORIAL 2 — MANUAL ESCALATION
% ================================================================
\begin{frame}
\frametitle{\fontsize{14}{17}\selectfont TUTORIAL 2: Manual Escalation of Suspicious Activity}

\small
\begin{block}{\fontsize{11}{13}\selectfont Topic: Frontline Staff Responsibilities}
\footnotesize
Frontline staff are often the first to notice suspicious activity. Even without a system alert, they must escalate concerns.
\end{block}

\vspace{0.03cm}

\stepbox{\footnotesize
\steptitle{Step 1:} Understand the scenario. A frontline staff receives a large, unusually timed cash deposit from a long-standing customer at a small bank branch. There is no immediate system alert. \\
\steptitle{Step 2:} Recognize the red flags. \\
\quad - Large deposit: Unusual for this customer \\
\quad - Unusually timed: Outside normal patterns \\
\quad - Cash: Hard to trace \\
\steptitle{Step 3:} Understand the correct action. The staff should \textbf{manually escalate a report to compliance}. \\
\steptitle{Step 4:} See why this is correct. The system did not generate an alert, but the staff has a \textbf{subjective concern}. Human observation is still valuable. \\
\steptitle{Step 5:} Apply the lesson. Frontline staff should escalate when they observe suspicious activity, even without a system alert.}

\vspace{0.08cm}
\tipbox{\footnotesize \textbf{CORE LESSON:} Frontline staff must escalate suspicious activity even when there is no system alert. Human observation is still a critical control.}
\end{frame}

% ================================================================
% SLIDE 4: TUTORIAL 3 — INVESTIGATION DOCUMENTATION
% ================================================================
\begin{frame}
\frametitle{\fontsize{14}{17}\selectfont TUTORIAL 3: Documentation Best Practices}

\small
\begin{block}{\fontsize{11}{13}\selectfont Topic: Suspicious Activity Investigations}
\footnotesize
Effective documentation is essential for suspicious activity investigations. It provides a record of the investigation and supports decision-making.
\end{block}

\vspace{0.03cm}

\stepbox{\footnotesize
\steptitle{Step 1:} Understand the importance of documentation. Investigations must be documented to: \\
\quad - Provide an audit trail \\
\quad - Support decision-making \\
\quad - Demonstrate compliance \\
\steptitle{Step 2:} Learn the first best practice. \textbf{Using structured templates for investigative notes} ensures consistency and completeness. \\
\steptitle{Step 3:} Learn the second best practice. \textbf{Adapting tools based on jurisdictional requirements} ensures compliance with local regulations. \\
\steptitle{Step 4:} See what is NOT the best practice. \textbf{Documenting findings only at the conclusion} is insufficient. Documentation should be ongoing throughout the investigation. \\
\steptitle{Step 5:} Apply the lesson. Use structured templates and adapt tools to jurisdictional requirements. Document throughout the investigation, not just at the end.}

\vspace{0.08cm}
\tipbox{\footnotesize \textbf{CORE LESSON:} Use structured templates for investigative notes and adapt tools to jurisdictional requirements. Document throughout, not just at the conclusion.}
\end{frame}

% ================================================================
% SLIDE 5: TUTORIAL 4 — BYPASSING APPROVAL PROCESSES
% ================================================================
\begin{frame}
\frametitle{\fontsize{14}{17}\selectfont TUTORIAL 4: Bypassing Approval Processes}

\small
\begin{block}{\fontsize{11}{13}\selectfont Topic: Trade Finance Controls}
\footnotesize
When employees bypass approval processes, it may indicate money laundering activity. Understanding the concern is essential.
\end{block}

\vspace{0.03cm}

\stepbox{\footnotesize
\steptitle{Step 1:} Understand the scenario. An employee repeatedly bypassed approval processes for low-dollar-value international transactions to a jurisdiction known for banking secrecy. \\
\steptitle{Step 2:} Recognize the compliance concern. This is \textbf{potential layering activity characteristic of money laundering}. \\
\steptitle{Step 3:} Understand why this is concerning. \\
\quad - Bypassing approvals: The employee is avoiding oversight \\
\quad - Low-dollar values: Trying to stay below thresholds \\
\quad - Banking secrecy jurisdiction: Trying to hide the activity \\
\steptitle{Step 4:} See the pattern. This combination suggests layering — moving money to hide its origin. \\
\steptitle{Step 5:} Apply the lesson. Repeated bypassing of approvals, especially to secrecy jurisdictions, is a red flag for money laundering.}

\vspace{0.08cm}
\tipbox{\footnotesize \textbf{CORE LESSON:} Bypassing approvals for transactions to secrecy jurisdictions suggests layering activity. This is a significant compliance concern.}
\end{frame}

% ================================================================
% SLIDE 6: TUTORIAL 5 — TIERED ONBOARDING FOR E-COMMERCE
% ================================================================
\begin{frame}
\frametitle{\fontsize{14}{17}\selectfont TUTORIAL 5: E-Commerce Risk Mitigation}

\small
\begin{block}{\fontsize{11}{13}\selectfont Topic: Money Laundering Risks in Nonbank Financial Services}
\footnotesize
E-commerce platforms face unique money laundering risks. Specific strategies help mitigate these risks.
\end{block}

\vspace{0.03cm}

\stepbox{\footnotesize
\steptitle{Step 1:} Understand the risks. E-commerce platforms can be used for money laundering through: \\
\quad - Fake transactions \\
\quad - Over/under invoicing \\
\quad - Layering through multiple sellers \\
\steptitle{Step 2:} Learn the first strategy. \textbf{Introducing tiered onboarding based on projected sales activity} helps manage risk by scaling due diligence. \\
\steptitle{Step 3:} Learn the second strategy. \textbf{Conducting sanctions screening on all sellers during onboarding} ensures that sanctioned individuals cannot use the platform. \\
\steptitle{Step 4:} See what is NOT the best strategy. \textbf{Implementing seller-rankings based on transaction volume} does not address money laundering risk — it may actually incentivize volume. \\
\steptitle{Step 5:} Apply the lesson. Mitigate e-commerce risk with tiered onboarding based on risk and sanctions screening.}

\vspace{0.08cm}
\tipbox{\footnotesize \textbf{CORE LESSON:} E-commerce platforms should use tiered onboarding based on projected sales activity and sanction all sellers.}
\end{frame}

% ================================================================
% SLIDE 7: TUTORIAL 6 — TOKENS AND DECENTRALIZED PLATFORMS
% ================================================================
\begin{frame}
\frametitle{\fontsize{14}{17}\selectfont TUTORIAL 6: Tokens and Decentralized Platforms}

\small
\begin{block}{\fontsize{11}{13}\selectfont Topic: Cryptocurrency and Money Laundering}
\footnotesize
Tokens on decentralized platforms present unique money laundering risks. Understanding these risks is essential.
\end{block}

\vspace{0.03cm}

\stepbox{\footnotesize
\steptitle{Step 1:} Understand the risk. \textbf{Tokens may be exchanged rapidly on decentralized platforms, obscuring source and destination of funds}. \\
\steptitle{Step 2:} Recognize why decentralized platforms are risky. \\
\quad - No central authority to monitor \\
\quad - Transactions are pseudonymous \\
\quad - Rapid exchanges make tracing difficult \\
\steptitle{Step 3:} Compare to centralized platforms. Centralized platforms have compliance controls and can be monitored. Decentralized platforms lack these controls. \\
\steptitle{Step 4:} Understand the laundering technique. Criminals can rapidly exchange tokens through multiple wallets, making it nearly impossible to trace the source. \\
\steptitle{Step 5:} Apply the lesson. Decentralized platforms and tokens present significant AML challenges due to the lack of central oversight.}

\vspace{0.08cm}
\tipbox{\footnotesize \textbf{CORE LESSON:} Tokens on decentralized platforms can be exchanged rapidly, obscuring source and destination. This makes them attractive for money laundering.}
\end{frame}

% ================================================================
% SLIDE 8: TUTORIAL 7 — RECOMMENDATION 18
% ================================================================
\begin{frame}
\frametitle{\fontsize{14}{17}\selectfont TUTORIAL 7: FATF Recommendation 18}

\small
\begin{block}{\fontsize{11}{13}\selectfont Topic: Global AFC Standards and Guidance}
\footnotesize
FATF Recommendation 18 addresses group-wide AML programs. Deficiencies in implementation are a common audit finding.
\end{block}

\vspace{0.03cm}

\stepbox{\footnotesize
\steptitle{Step 1:} Understand what Recommendation 18 covers. It requires financial groups to implement \textbf{group-wide AML programs that apply to all branches and subsidiaries}. \\
\steptitle{Step 2:} Recognize the focus of an audit. Auditors reviewing Recommendation 18 implementation would focus on \textbf{whether group-wide AML programs apply to all branches and subsidiaries}. \\
\steptitle{Step 3:} Understand why this matters. Without group-wide programs: \\
\quad - Some branches may have weaker controls \\
\quad - Criminals can exploit the weakest link \\
\quad - The group is exposed to risk \\
\steptitle{Step 4:} See what is NOT the focus. \textbf{The training curriculum for non-AML business units} is not the focus of Recommendation 18. \\
\steptitle{Step 5:} Apply the lesson. Recommendation 18 requires group-wide AML programs across all branches and subsidiaries.}

\vspace{0.08cm}
\tipbox{\footnotesize \textbf{CORE LESSON:} FATF Recommendation 18 requires group-wide AML programs that apply to all branches and subsidiaries. Auditors focus on whether programs are truly group-wide.}
\end{frame}

% ================================================================
% SLIDE 9: TUTORIAL 8 — FSRB ROLE
% ================================================================
\begin{frame}
\frametitle{\fontsize{14}{17}\selectfont TUTORIAL 8: FATF-Style Regional Bodies}

\small
\begin{block}{\fontsize{11}{13}\selectfont Topic: Global AFC Standards}
\footnotesize
FATF-Style Regional Bodies (FSRBs) play a key role in implementing and promoting FATF standards.
\end{block}

\vspace{0.03cm}

\stepbox{\footnotesize
\steptitle{Step 1:} Understand what FSRBs are. They are regional organizations that \textbf{help implement and promote FATF standards by evaluating compliance in their regions}. \\
\steptitle{Step 2:} Recognize their role. FSRBs: \\
\quad - Conduct mutual evaluations in their region \\
\quad - Provide technical assistance \\
\quad - Promote FATF standards \\
\steptitle{Step 3:} See the relationship to FATF. FATF sets global standards. FSRBs implement and enforce them at the regional level. \\
\steptitle{Step 4:} Understand the value. FSRBs bring FATF standards closer to national implementation. They provide local expertise and context. \\
\steptitle{Step 5:} Apply the lesson. FSRBs are essential for the effective implementation of FATF standards at the regional level.}

\vspace{0.08cm}
\tipbox{\footnotesize \textbf{CORE LESSON:} FSRBs help implement and promote FATF standards by evaluating compliance in their regions. They are the regional enforcers of global standards.}
\end{frame}

% ================================================================
% SLIDE 10: TUTORIAL 9 — IOSCO OBJECTIVES
% ================================================================
\begin{frame}
\frametitle{\fontsize{14}{17}\selectfont TUTORIAL 9: IOSCO Core Objectives}

\small
\begin{block}{\fontsize{11}{13}\selectfont Topic: International Organization of Securities Commissions}
\footnotesize
IOSCO's core objectives support anti-financial crime efforts in securities markets. Understanding these objectives is essential.
\end{block}

\vspace{0.03cm}

\stepbox{\footnotesize
\steptitle{Step 1:} Understand IOSCO's role. IOSCO sets global standards for securities regulation. \\
\steptitle{Step 2:} Learn the objectives that support AFC. Two objectives are: \\
\quad (1) \textbf{Enhancing investor protection through cooperation and enforcement} \\
\quad (2) \textbf{Promoting financial stability by managing systemic risks and facilitating international information exchange} \\
\steptitle{Step 3:} Understand why these support AFC. \\
\quad - Investor protection includes preventing financial crime \\
\quad - Cooperation and enforcement enable cross-border action \\
\quad - Financial stability includes preventing money laundering \\
\quad - Information exchange supports investigations \\
\steptitle{Step 4:} See what is NOT correct. \textbf{Requiring uniform compliance department structures} is not an IOSCO objective. \\
\steptitle{Step 5:} Apply the lesson. IOSCO supports AFC through investor protection and financial stability objectives.}

\vspace{0.08cm}
\tipbox{\footnotesize \textbf{CORE LESSON:} IOSCO supports anti-financial crime through investor protection (cooperation and enforcement) and financial stability (risk management and information exchange).}
\end{frame}

% ================================================================
% SLIDE 11: TUTORIAL 10 — IMF EVALUATIONS
% ================================================================
\begin{frame}
\frametitle{\fontsize{14}{17}\selectfont TUTORIAL 10: IMF Financial System Evaluations}

\small
\begin{block}{\fontsize{11}{13}\selectfont Topic: IMF Assessments}
\footnotesize
The IMF evaluates both economic measures and financial system vulnerabilities. Economic improvements can mask financial crime risks.
\end{block}

\vspace{0.03cm}

\stepbox{\footnotesize
\steptitle{Step 1:} Understand the scenario. A nation achieved strong GDP growth and reduced inflation after implementing IMF-recommended fiscal reforms. However, the IMF still expresses significant concerns. \\
\steptitle{Step 2:} Understand the explanation. \textbf{Improvements in economic measures can mask financial crime vulnerabilities}. \\
\steptitle{Step 3:} Recognize why this happens. \\
\quad - GDP growth can hide illicit flows \\
\quad - Reduced inflation does not mean reduced money laundering \\
\quad - Economic reforms do not automatically address AML \\
\steptitle{Step 4:} Understand the IMF's approach. The IMF looks at both economic and financial system health. Financial crime vulnerabilities can exist even with strong economic indicators. \\
\steptitle{Step 5:} Apply the lesson. Strong economic measures do not guarantee financial system integrity. Financial crime risks may still exist.}

\vspace{0.08cm}
\tipbox{\footnotesize \textbf{CORE LESSON:} Improvements in economic measures can mask financial crime vulnerabilities. Strong GDP growth does not mean strong AML controls.}
\end{frame}

% ================================================================
% SLIDE 12: TUTORIAL 11 — FALSE POSITIVES AND SCREENING FATIGUE
% ================================================================
\begin{frame}
\frametitle{\fontsize{14}{17}\selectfont TUTORIAL 11: False Positives and Screening Fatigue}

\small
\begin{block}{\fontsize{11}{13}\selectfont Topic: AML Program Controls}
\footnotesize
High numbers of false positives in payment screening systems create significant problems. Understanding how to address this is essential.
\end{block}

\vspace{0.03cm}

\stepbox{\footnotesize
\steptitle{Step 1:} Understand the scenario. A payment screening system is returning a high number of false positives. \\
\steptitle{Step 2:} Recognize the impact. \textbf{It may cause screening fatigue} — analysts become desensitized to alerts and may miss real risks. \\
\steptitle{Step 3:} Understand the solution. \textbf{Tuning the scenario logic or data inputs} will help reduce false positives. \\
\steptitle{Step 4:} See what is NOT the solution. \textbf{Reducing list matching criteria to minimize alerts} is wrong — it would miss real matches. \\
\steptitle{Step 5:} Apply the lesson. Address false positives by tuning the system, not by reducing criteria.}

\vspace{0.08cm}
\tipbox{\footnotesize \textbf{CORE LESSON:} High false positives cause screening fatigue. Tune the scenario logic or data inputs rather than reducing matching criteria.}
\end{frame}

% ================================================================
% SLIDE 13: TUTORIAL 12 — AUTOMATED ADVERSE MEDIA CHECKS
% ================================================================
\begin{frame}
\frametitle{\fontsize{14}{17}\selectfont TUTORIAL 12: Automated Adverse Media Checks}

\small
\begin{block}{\fontsize{11}{13}\selectfont Topic: AML Technology}
\footnotesize
Automated systems for adverse media checks provide significant advantages over manual systems.
\end{block}

\vspace{0.03cm}

\stepbox{\footnotesize
\steptitle{Step 1:} Understand what adverse media checks are. They involve searching for negative news about customers or entities. \\
\steptitle{Step 2:} Recognize the challenge of manual checks. Manual checks are: \\
\quad - Time-consuming \\
\quad - Limited in scope \\
\quad - Prone to errors \\
\steptitle{Step 3:} Understand the advantage of automation. Automated systems \textbf{improve efficiency in identifying negative news across multiple sources}. \\
\steptitle{Step 4:} See the benefits of automation. \\
\quad - Covers more sources \\
\quad - Faster processing \\
\quad - More consistent \\
\steptitle{Step 5:} Apply the lesson. Automation improves efficiency and coverage in adverse media checks.}

\vspace{0.08cm}
\tipbox{\footnotesize \textbf{CORE LESSON:} Automated adverse media checks improve efficiency in identifying negative news across multiple sources, unlike manual checks.}
\end{frame}

% ================================================================
% SLIDE 14: TUTORIAL 13 — AI-BASED SYSTEM INCONSISTENCY
% ================================================================
\begin{frame}
\frametitle{\fontsize{14}{17}\selectfont TUTORIAL 13: AI-Based Screening Inconsistency}

\small
\begin{block}{\fontsize{11}{13}\selectfont Topic: Technology for AFC Compliance}
\footnotesize
When AI-based systems reject transactions inconsistently across business units, corrective action is needed.
\end{block}

\vspace{0.03cm}

\stepbox{\footnotesize
\steptitle{Step 1:} Understand the scenario. A newly implemented AI-based screening system is rejecting transactions inconsistently across business units. \\
\steptitle{Step 2:} Recognize the problem. Inconsistency means the system is not applying rules uniformly. This creates: \\
\quad - Unfair treatment of customers \\
\quad - Inconsistent compliance \\
\quad - Regulatory risk \\
\steptitle{Step 3:} Understand the correct first step. \textbf{Engage the governance committee to assess model performance}. \\
\steptitle{Step 4:} See why this is the right approach. The governance committee: \\
\quad - Has oversight responsibility \\
\quad - Can evaluate the model \\
\quad - Can authorize changes \\
\steptitle{Step 5:} Apply the lesson. When AI systems perform inconsistently, engage the governance committee for assessment.}

\vspace{0.08cm}
\tipbox{\footnotesize \textbf{CORE LESSON:} When AI screening systems are inconsistent, the governance committee should assess model performance.}
\end{frame}

% ================================================================
% SLIDE 15: TUTORIAL 14 — DEVICE INTELLIGENCE TOOLS
% ================================================================
\begin{frame}
\frametitle{\fontsize{14}{17}\selectfont TUTORIAL 14: Device Intelligence Tools}

\small
\begin{block}{\fontsize{11}{13}\selectfont Topic: Technology for AFC Compliance}
\footnotesize
Device intelligence tools capture specific elements to enhance detection of online financial crimes.
\end{block}

\vspace{0.03cm}

\stepbox{\footnotesize
\steptitle{Step 1:} Understand what device intelligence tools do. They capture information about devices used to access financial systems. \\
\steptitle{Step 2:} Learn the elements captured. Three key elements are: \\
\quad (1) \textbf{Browser type and screen resolution} \\
\quad (2) \textbf{IP address and geolocation} \\
\quad (3) \textbf{Operating system and installed plugins} \\
\steptitle{Step 3:} Understand why these matter. \\
\quad - Browser/screen: Identifies the device type \\
\quad - IP/geolocation: Identifies the location \\
\quad - OS/plugins: Identifies the environment \\
\steptitle{Step 4:} See what is NOT typically captured. \textbf{Social media accounts and browsing history} are not typically captured by device intelligence tools. \\
\steptitle{Step 5:} Apply the lesson. Device intelligence captures device-specific information, not social media or browsing history.}

\vspace{0.08cm}
\tipbox{\footnotesize \textbf{CORE LESSON:} Device intelligence tools capture browser type, IP address/geolocation, and operating system/plugins. They do not capture social media accounts.}
\end{frame}

% ================================================================
% SLIDE 16: TUTORIAL 15 — WIRES THROUGH SMALL FOREIGN BANKS
% ================================================================
\begin{frame}
\frametitle{\fontsize{14}{17}\selectfont TUTORIAL 15: Wire Transfers Through Small Foreign Banks}

\small
\begin{block}{\fontsize{11}{13}\selectfont Topic: Transaction Monitoring}
\footnotesize
Multiple wire transfers through small foreign banks without clear business need require investigation.
\end{block}

\vspace{0.03cm}

\stepbox{\footnotesize
\steptitle{Step 1:} Understand the scenario. Multiple wire transfers from a single client are being routed through small foreign banks without a clear business need. \\
\steptitle{Step 2:} Recognize the red flags. \\
\quad - Small foreign banks: May have weaker controls \\
\quad - No clear business need: Unusual for this client \\
\quad - Multiple transfers: Suggests structuring \\
\steptitle{Step 3:} Understand the correct action. \textbf{Investigate further by requesting justification and reviewing wire documentation}. \\
\steptitle{Step 4:} See why this is the right approach. You need to understand: \\
\quad - Why these banks are used \\
\quad - The purpose of the transfers \\
\quad - Whether the documentation supports the activity \\
\steptitle{Step 5:} Apply the lesson. When clients use unusual banking channels without business need, investigate further.}

\vspace{0.08cm}
\tipbox{\footnotesize \textbf{CORE LESSON:} Wires through small foreign banks without business need require investigation. Request justification and review documentation.}
\end{frame}

% ================================================================
% SLIDE 17: TUTORIAL 16 — TRADE-BASED MONEY LAUNDERING
% ================================================================
\begin{frame}
\frametitle{\fontsize{14}{17}\selectfont TUTORIAL 16: Trade-Based Money Laundering Controls}

\small
\begin{block}{\fontsize{11}{13}\selectfont Topic: Trade Finance Risk Reduction}
\footnotesize
Specific operational and product changes can effectively reduce trade-based money laundering exposure.
\end{block}

\vspace{0.03cm}

\stepbox{\footnotesize
\steptitle{Step 1:} Understand the need for controls. Trade finance is vulnerable to TBML through invoice fraud, over/under invoicing, and misrepresentation. \\
\steptitle{Step 2:} Learn the first control. \textbf{Require digitized bills of lading for all high-value shipments} — digitization makes manipulation harder. \\
\steptitle{Step 3:} Learn the second control. \textbf{Implement dual-person approval for customer invoice verification above threshold} — segregation of duties prevents fraud. \\
\steptitle{Step 4:} Learn the third control. \textbf{Increase transaction look-back periods in system audits} — longer look-back catches patterns. \\
\steptitle{Step 5:} Apply the lesson. TBML exposure is reduced through digitization, dual approvals, and longer look-back periods.}

\vspace{0.08cm}
\tipbox{\footnotesize \textbf{CORE LESSON:} Reduce TBML exposure with digitized bills of lading, dual-person approval for invoice verification, and increased transaction look-back periods.}
\end{frame}

% ================================================================
% SLIDE 18: TUTORIAL 17 — SYNTHETIC IDENTITY FRAUD
% ================================================================
\begin{frame}
\frametitle{\fontsize{14}{17}\selectfont TUTORIAL 17: Synthetic Identity Fraud Response}

\small
\begin{block}{\fontsize{11}{13}\selectfont Topic: Product Risk Management}
\footnotesize
When synthetic identity fraud is linked to a new product, corrective action is needed.
\end{block}

\vspace{0.03cm}

\stepbox{\footnotesize
\steptitle{Step 1:} Understand the scenario. A new credit product targeting low-income customers is experiencing a trend of synthetic identity fraud. \\
\steptitle{Step 2:} Recognize the problem. Synthetic identity fraud involves creating fake identities to obtain credit. This is a serious risk for new products. \\
\steptitle{Step 3:} Understand the correct action. \textbf{Suspend the product while updating the risk assessment and onboarding controls}. \\
\steptitle{Step 4:} See why this is the right approach. \\
\quad - Suspension stops further losses \\
\quad - Risk assessment identifies the vulnerabilities \\
\quad - Updated controls prevent recurrence \\
\steptitle{Step 5:} Apply the lesson. When fraud is detected in a product, suspend it, update the risk assessment, and strengthen controls.}

\vspace{0.08cm}
\tipbox{\footnotesize \textbf{CORE LESSON:} When synthetic identity fraud is detected, suspend the product, update the risk assessment, and strengthen onboarding controls.}
\end{frame}

% ================================================================
% SLIDE 19: TUTORIAL 18 — PRIVATE BANKING HIGH-RISK FEATURES
% ================================================================
\begin{frame}
\frametitle{\fontsize{14}{17}\selectfont TUTORIAL 18: Private Banking Risk Features}

\small
\begin{block}{\fontsize{11}{13}\selectfont Topic: Private Banking Risks}
\footnotesize
Private banking has specific features that create high risk for money laundering.
\end{block}

\vspace{0.03cm}

\stepbox{\footnotesize
\steptitle{Step 1:} Understand the context. Private banking involves high-net-worth individuals and complex services. \\
\steptitle{Step 2:} Learn the high-risk features. Three key features are: \\
\quad (1) \textbf{Offering numbered accounts} — anonymizes ownership \\
\quad (2) \textbf{Dedicated relationship service with transaction discretion} — creates intimacy and opportunity \\
\quad (3) \textbf{Minimum assets requirement} — high-value clients with complex needs \\
\steptitle{Step 3:} Understand why these create risk. \\
\quad - Numbered accounts: Hide identity \\
\quad - Relationship discretion: Reduce oversight \\
\quad - High assets: More money to launder \\
\steptitle{Step 4:} Apply the lesson. Private banking risk is highest when accounts are numbered, relationships are discretionary, and assets are high.}

\vspace{0.08cm}
\tipbox{\footnotesize \textbf{CORE LESSON:} Private banking high-risk features include numbered accounts, dedicated relationship service with discretion, and minimum assets requirement.}
\end{frame}

% ================================================================
% SLIDE 20: TUTORIAL 19 — EGMONT GROUP BENEFITS
% ================================================================
\begin{frame}
\frametitle{\fontsize{14}{17}\selectfont TUTORIAL 19: Egmont Group Benefits}

\small
\begin{block}{\fontsize{11}{13}\selectfont Topic: Financial Intelligence Units}
\footnotesize
Egmont Group membership provides significant benefits for national Financial Intelligence Units.
\end{block}

\vspace{0.03cm}

\stepbox{\footnotesize
\steptitle{Step 1:} Understand what the Egmont Group is. It is a global network of Financial Intelligence Units (FIUs). \\
\steptitle{Step 2:} Learn the benefits of membership. Two key benefits are: \\
\quad (1) \textbf{Access to a secure online portal for communication with peer FIUs} \\
\quad (2) \textbf{Increased capacity through shared typology projects and analyses} \\
\steptitle{Step 3:} Understand why these matter. \\
\quad - Secure portal: Enables confidential information sharing \\
\quad - Shared typologies: Helps identify emerging threats \\
\steptitle{Step 4:} See what is NOT a benefit. \textbf{Automatic inclusion in FATF mutual evaluations} is not a benefit of Egmont membership. \\
\steptitle{Step 5:} Apply the lesson. Egmont membership provides secure communication and shared analysis capabilities.}

\vspace{0.08cm}
\tipbox{\footnotesize \textbf{CORE LESSON:} Egmont Group membership provides secure communication with peer FIUs and increased capacity through shared typologies.}
\end{frame}

% ================================================================
% SLIDE 21: TUTORIAL 20 — FATF MUTUAL EVALUATION RESPONSE
% ================================================================
\begin{frame}
\frametitle{\fontsize{14}{17}\selectfont TUTORIAL 20: Responding to FATF Mutual Evaluations}

\small
\begin{block}{\fontsize{11}{13}\selectfont Topic: FATF Mutual Evaluations}
\footnotesize
When a FATF mutual evaluation identifies deficiencies, the country must take corrective action.
\end{block}

\vspace{0.03cm}

\stepbox{\footnotesize
\steptitle{Step 1:} Understand the scenario. A central bank head is responding to a negative FATF mutual evaluation that highlighted deficiencies in beneficial ownership transparency and transaction reporting. \\
\steptitle{Step 2:} Recognize the deficiencies. \\
\quad - Beneficial ownership transparency: Who really owns companies? \\
\quad - Transaction reporting: Are suspicious transactions being reported? \\
\steptitle{Step 3:} Understand the best response. \textbf{Increase AML training for compliance officers and mandate tighter oversight of beneficial ownership records}. \\
\steptitle{Step 4:} See why this works. \\
\quad - Training: Improves skills and knowledge \\
\quad - Tighter oversight: Ensures compliance \\
\steptitle{Step 5:} Apply the lesson. Respond to FATF deficiencies with training and stronger oversight.}

\vspace{0.08cm}
\tipbox{\footnotesize \textbf{CORE LESSON:} Respond to FATF mutual evaluation deficiencies with increased AML training and tighter oversight of beneficial ownership records.}
\end{frame}

% ================================================================
% SLIDE 22: TUTORIAL 21 — EDD FOR UNRELATED CUSTOMERS
% ================================================================
\begin{frame}
\frametitle{\fontsize{14}{17}\selectfont TUTORIAL 21: Enhanced Due Diligence for Common Beneficial Owner}

\small
\begin{block}{\fontsize{11}{13}\selectfont Topic: Customer Due Diligence}
\footnotesize
When unrelated customers share a beneficial owner in a high-risk jurisdiction, enhanced due diligence is required.
\end{block}

\vspace{0.03cm}

\stepbox{\footnotesize
\steptitle{Step 1:} Understand the scenario. Three apparently unrelated customers share the same beneficial owner. The owner is in a high-risk jurisdiction and was not mentioned in onboarding documentation. \\
\steptitle{Step 2:} Recognize the red flags. \\
\quad - Shared beneficial owner: They are related \\
\quad - High-risk jurisdiction: The owner is in a risk area \\
\quad - Not mentioned before: Information was hidden \\
\steptitle{Step 3:} Understand the correct action. \textbf{Conduct enhanced due diligence to determine the legitimacy of the relationships}. \\
\steptitle{Step 4:} See why this is the right approach. EDD is needed to: \\
\quad - Understand the relationships \\
\quad - Verify the information \\
\quad - Assess the risk \\
\steptitle{Step 5:} Apply the lesson. When customers share a beneficial owner in a high-risk jurisdiction, conduct enhanced due diligence.}

\vspace{0.08cm}
\tipbox{\footnotesize \textbf{CORE LESSON:} When unrelated customers share a beneficial owner in a high-risk jurisdiction, conduct enhanced due diligence to determine legitimacy.}
\end{frame}

% ================================================================
% SLIDE 23: TUTORIAL 22 — CONCENTRATION ACCOUNT USES
% ================================================================
\begin{frame}
\frametitle{\fontsize{14}{17}\selectfont TUTORIAL 22: Concentration Account Uses}

\small
\begin{block}{\fontsize{11}{13}\selectfont Topic: Concentration Accounts}
\footnotesize
Concentration accounts have legitimate uses within financial institutions. Understanding these uses is essential.
\end{block}

\vspace{0.03cm}

\stepbox{\footnotesize
\steptitle{Step 1:} Understand what concentration accounts are. They are accounts used to aggregate and move funds internally. \\
\steptitle{Step 2:} Learn the legitimate uses. Two typical uses are: \\
\quad (1) \textbf{Streamlining inter-branch cash movements within the bank} \\
\quad (2) \textbf{Facilitating back-office functions like suspense clearing} \\
\steptitle{Step 3:} Understand why these are legitimate. \\
\quad - Inter-branch cash movements: Efficient internal transfers \\
\quad - Suspense clearing: Temporary holding for processing \\
\steptitle{Step 4:} See what is NOT a legitimate use. \textbf{Aggregating transactions from multiple unrelated customers} creates traceability issues. \\
\steptitle{Step 5:} Apply the lesson. Concentration accounts are legitimately used for internal cash movements and back-office processing.}

\vspace{0.08cm}
\tipbox{\footnotesize \textbf{CORE LESSON:} Concentration accounts are legitimately used for inter-branch cash movements and back-office functions like suspense clearing.}
\end{frame}

% ================================================================
% SLIDE 24: TUTORIAL 23 — PEP RISK MANAGEMENT
% ================================================================
\begin{frame}
\frametitle{\fontsize{14}{17}\selectfont TUTORIAL 23: PEP Risk Management}

\small
\begin{block}{\fontsize{11}{13}\selectfont Topic: Politically Exposed Persons}
\footnotesize
Managing PEP-related money laundering risks requires specific measures.
\end{block}

\vspace{0.03cm}

\stepbox{\footnotesize
\steptitle{Step 1:} Understand the PEP risk. PEPs are more vulnerable to corruption and bribery. \\
\steptitle{Step 2:} Learn the effective measures. Three key measures are: \\
\quad (1) \textbf{Establish a system that auto-detects public official occupations} \\
\quad (2) \textbf{Capture information about source of funds and source of wealth at onboarding} \\
\quad (3) \textbf{Conduct ongoing media monitoring for adverse news on clients} \\
\steptitle{Step 3:} Understand why these are effective. \\
\quad - Auto-detection: Identifies PEPs early \\
\quad - Source of funds/wealth: Verifies legitimacy \\
\quad - Media monitoring: Identifies emerging risks \\
\steptitle{Step 4:} Apply the lesson. Manage PEP risk through detection, verification, and ongoing monitoring.}

\vspace{0.08cm}
\tipbox{\footnotesize \textbf{CORE LESSON:} Manage PEP risks with auto-detection of public officials, source of funds/wealth verification, and ongoing media monitoring.}
\end{frame}

% ================================================================
% SLIDE 25: TUTORIAL 24 — PREPAID CARD RISKS
% ================================================================
\begin{frame}
\frametitle{\fontsize{14}{17}\selectfont TUTORIAL 24: Prepaid Card Risks}

\small
\begin{block}{\fontsize{11}{13}\selectfont Topic: Payment Product Risks}
\footnotesize
Prepaid cards are particularly vulnerable to financial crime risks.
\end{block}

\vspace{0.03cm}

\stepbox{\footnotesize
\steptitle{Step 1:} Understand what prepaid cards are. They are payment cards that are pre-loaded with funds. \\
\steptitle{Step 2:} Recognize the vulnerability. \textbf{They can be used anonymously, especially when reloadable or unregistered}. \\
\steptitle{Step 3:} Understand why this creates risk. \\
\quad - Anonymous: Cannot trace to a person \\
\quad - Reloadable: Can add more funds \\
\quad - Unregistered: No identity verification \\
\steptitle{Step 4:} See how criminals exploit this. \\
\quad - Money laundering: Load and spend anonymously \\
\quad - Terrorist financing: Move funds without detection \\
\steptitle{Step 5:} Apply the lesson. Prepaid cards, especially reloadable and unregistered, are high-risk for financial crime.}

\vspace{0.08cm}
\tipbox{\footnotesize \textbf{CORE LESSON:} Prepaid cards are vulnerable because they can be used anonymously, especially when reloadable or unregistered.}
\end{frame}

% ================================================================
% SLIDE 26: TUTORIAL 25 — SANCTIONS SCREENING LISTS
% ================================================================
\begin{frame}
\frametitle{\fontsize{14}{17}\selectfont TUTORIAL 25: Sanctions Screening Lists}

\small
\begin{block}{\fontsize{11}{13}\selectfont Topic: Sanctions Controls}
\footnotesize
External sanctions lists are essential for effective sanctions screening controls.
\end{block}

\vspace{0.03cm}

\stepbox{\footnotesize
\steptitle{Step 1:} Understand what sanctions lists are. They identify individuals and entities subject to sanctions. \\
\steptitle{Step 2:} Learn the essential external sources. Two key sources are: \\
\quad (1) \textbf{FATF blacklists and UN Security Council lists} \\
\quad (2) \textbf{State Department embargo registries} \\
\steptitle{Step 3:} Understand why these are essential. \\
\quad - FATF/UN: Global standards \\
\quad - State Department: US-specific requirements \\
\steptitle{Step 4:} See what is NOT an external source. \textbf{Watchlists from the firm's previous investigations} are internal, not external. \\
\steptitle{Step 5:} Apply the lesson. Use authoritative external sources for sanctions screening.}

\vspace{0.08cm}
\tipbox{\footnotesize \textbf{CORE LESSON:} Essential sanctions screening lists include FATF blacklists, UN Security Council lists, and State Department embargo registries.}
\end{frame}

% ================================================================
% SLIDE 27: TUTORIAL 26 — INTERNAL VS EXTERNAL SOURCES
% ================================================================
\begin{frame}
\frametitle{\fontsize{14}{17}\selectfont TUTORIAL 26: Internal vs External Data Sources}

\small
\begin{block}{\fontsize{11}{13}\selectfont Topic: Data Sources for AFC}
\footnotesize
Internal and external data sources serve different purposes in AFC data extraction.
\end{block}

\vspace{0.03cm}

\stepbox{\footnotesize
\steptitle{Step 1:} Understand the distinction. Internal and external sources provide different types of information. \\
\steptitle{Step 2:} Learn the characteristics. \\
\quad - \textbf{Internal sources reflect observed customer behavior or system-held values} \\
\quad - \textbf{External sources include sanction or adverse media lists from agencies} \\
\steptitle{Step 3:} Understand why this distinction matters. \\
\quad - Internal: What the institution knows about the customer \\
\quad - External: What the world knows about the customer \\
\steptitle{Step 4:} Apply the lesson. Use both internal and external sources for a complete picture.}

\vspace{0.08cm}
\tipbox{\footnotesize \textbf{CORE LESSON:} Internal sources reflect customer behavior and system data. External sources include sanctions and adverse media lists.}
\end{frame}

% ================================================================
% SLIDE 28: TUTORIAL 27 — DATA LINEAGE
% ================================================================
\begin{frame}
\frametitle{\fontsize{14}{17}\selectfont TUTORIAL 27: Data Lineage for Transparency}

\small
\begin{block}{\fontsize{11}{13}\selectfont Topic: Data Governance — Lineage}
\footnotesize
Data lineage supports transparency and auditability in compliance systems.
\end{block}

\vspace{0.03cm}

\stepbox{\footnotesize
\steptitle{Step 1:} Understand what data lineage is. It is the history of data — where it came from and how it changed. \\
\steptitle{Step 2:} Learn the elements that support transparency. Two key elements are: \\
\quad (1) \textbf{Maintaining history of data transformations} \\
\quad (2) \textbf{Capturing source of each data field} \\
\steptitle{Step 3:} Understand why these support transparency. \\
\quad - Transformation history: Shows how data was modified \\
\quad - Source capture: Shows where data originated \\
\steptitle{Step 4:} Apply the lesson. Data lineage requires documenting data transformations and sources.}

\vspace{0.08cm}
\tipbox{\footnotesize \textbf{CORE LESSON:} Data lineage for transparency requires maintaining history of data transformations and capturing the source of each data field.}
\end{frame}

% ================================================================
% SLIDE 29: TUTORIAL 28 — BASEL INSTITUTE SERVICES
% ================================================================
\begin{frame}
\frametitle{\fontsize{14}{17}\selectfont TUTORIAL 28: Basel Institute Services}

\small
\begin{block}{\fontsize{11}{13}\selectfont Topic: Anti-Corruption Support}
\footnotesize
The Basel Institute provides specific services to support anti-corruption reforms.
\end{block}

\vspace{0.03cm}

\stepbox{\footnotesize
\steptitle{Step 1:} Understand the scenario. An NGO is developing anti-corruption mechanisms in a country with weak legal infrastructure and endemic fraud. \\
\steptitle{Step 2:} Recognize what is needed. The country needs both \textbf{legal advisory services} and \textbf{capacity building}. \\
\steptitle{Step 3:} Understand the correct service combination. \textbf{ICAR legal advisory services combined with capacity building programs} provide the most strategic support. \\
\steptitle{Step 4:} See why this works. \\
\quad - ICAR: Provides legal expertise \\
\quad - Capacity building: Builds local skills \\
\steptitle{Step 5:} Apply the lesson. Weak legal infrastructure requires both advisory services and capacity building.}

\vspace{0.08cm}
\tipbox{\footnotesize \textbf{CORE LESSON:} The Basel Institute's ICAR legal advisory services combined with capacity building programs provide strategic anti-corruption support.}
\end{frame}

% ================================================================
% SLIDE 30: TUTORIAL 29 — WORLD BANK/IMF EXPECTATIONS
% ================================================================
\begin{frame}
\frametitle{\fontsize{14}{17}\selectfont TUTORIAL 29: World Bank and IMF Expectations}

\small
\begin{block}{\fontsize{11}{13}\selectfont Topic: International Financial Institutions}
\footnotesize
The World Bank and IMF expect countries receiving their support to meet certain standards.
\end{block}

\vspace{0.03cm}

\stepbox{\footnotesize
\steptitle{Step 1:} Understand the scenario. A developing country receives significant funding from the World Bank. Internal audits reveal ongoing weaknesses in terrorism financing controls and legislative gaps. \\
\steptitle{Step 2:} Recognize the expectation. The country should \textbf{reform its AML/CFT frameworks and align them with international expectations}. \\
\steptitle{Step 3:} Understand why this is expected. \\
\quad - World Bank funding comes with conditions \\
\quad - International expectations include AML/CFT \\
\quad - Weak controls undermine development \\
\steptitle{Step 4:} Apply the lesson. Countries receiving international support must align with AML/CFT standards.}

\vspace{0.08cm}
\tipbox{\footnotesize \textbf{CORE LESSON:} Countries receiving World Bank/IMF support are expected to reform AML/CFT frameworks and align with international standards.}
\end{frame}

% ================================================================
% SLIDE 31: TUTORIAL 30 — G-20 ACWG INITIATIVES
% ================================================================
\begin{frame}
\frametitle{\fontsize{14}{17}\selectfont TUTORIAL 30: G-20 Anti-Corruption Working Group}

\small
\begin{block}{\fontsize{11}{13}\selectfont Topic: G-20 Anti-Corruption Efforts}
\footnotesize
The G-20 Anti-Corruption Working Group addresses key anti-corruption challenges.
\end{block}

\vspace{0.03cm}

\stepbox{\footnotesize
\steptitle{Step 1:} Understand the G-20 ACWG's role. It coordinates anti-corruption efforts among G-20 member states. \\
\steptitle{Step 2:} Understand the initiative that addresses asset recovery. \textbf{Increasing efficiency of asset recovery measures} most directly addresses the challenge of recovering illicitly obtained funds across jurisdictions. \\
\steptitle{Step 3:} Recognize why this matters. \\
\quad - Asset recovery is difficult across borders \\
\quad - Many countries lack effective mechanisms \\
\quad - Recovered funds can support development \\
\steptitle{Step 4:} Apply the lesson. Asset recovery efficiency is a key focus of international anti-corruption efforts.}

\vspace{0.08cm}
\tipbox{\footnotesize \textbf{CORE LESSON:} The G-20 ACWG addresses asset recovery challenges by focusing on increasing efficiency of asset recovery measures.}
\end{frame}

% ================================================================
% SLIDE 32: TUTORIAL 31 — PEP ESCALATION PROTOCOLS
% ================================================================
\begin{frame}
\frametitle{\fontsize{14}{17}\selectfont TUTORIAL 31: PEP Escalation Protocols}

\small
\begin{block}{\fontsize{11}{13}\selectfont Topic: Politically Exposed Persons}
\footnotesize
Frontline staff must know when to escalate alerts involving PEPs. BCBS guidance provides direction.
\end{block}

\vspace{0.03cm}

\stepbox{\footnotesize
\steptitle{Step 1:} Understand the scenario. Frontline staff are unsure when to escalate alerts involving PEPs. \\
\steptitle{Step 2:} Recognize the problem. Uncertainty leads to: \\
\quad - Missed escalations \\
\quad - Inconsistent treatment \\
\quad - Regulatory risk \\
\steptitle{Step 3:} Understand the corrective action. \textbf{Reinforce risk-based escalation protocols and provide KYC training}. \\
\steptitle{Step 4:} See why this works. \\
\quad - Protocols: Clear guidance \\
\quad - Training: Builds competence \\
\steptitle{Step 5:} Apply the lesson. BCBS guidance emphasizes clear escalation protocols and training for PEPs.}

\vspace{0.08cm}
\tipbox{\footnotesize \textbf{CORE LESSON:} Reinforce risk-based escalation protocols and provide KYC training to ensure staff know when to escalate PEP alerts.}
\end{frame}

% ================================================================
% SLIDE 33: TUTORIAL 32 — EU AML DIRECTIVE ALIGNMENT
% ================================================================
\begin{frame}
\frametitle{\fontsize{14}{17}\selectfont TUTORIAL 32: EU AML Directive Alignment}

\small
\begin{block}{\fontsize{11}{13}\selectfont Topic: EU AML Regulations}
\footnotesize
European banks must align with the latest EU AML Directives. Historical context is important.
\end{block}

\vspace{0.03cm}

\stepbox{\footnotesize
\steptitle{Step 1:} Understand the context. The EU has a history of AML regulation, with significant updates after the 2008 financial crisis. \\
\steptitle{Step 2:} Recognize the best strategy. \textbf{Incorporate recommendations from AML regulation updates post-2008 financial crisis}. \\
\steptitle{Step 3:} Understand why this is the best approach. \\
\quad - Post-2008 updates addressed crisis-related weaknesses \\
\quad - They represent current best practices \\
\quad - They align with international standards \\
\steptitle{Step 4:} Apply the lesson. Align with the latest AML Directive by incorporating post-2008 crisis updates.}

\vspace{0.08cm}
\tipbox{\footnotesize \textbf{CORE LESSON:} Align with EU AML Directives by incorporating recommendations from post-2008 financial crisis AML updates.}
\end{frame}

% ================================================================
% SLIDE 34: TUTORIAL 33 — SUSPICIOUS KYC ACCESS
% ================================================================
\begin{frame}
\frametitle{\fontsize{14}{17}\selectfont TUTORIAL 33: Investigating Suspicious KYC Access}

\small
\begin{block}{\fontsize{11}{13}\selectfont Topic: Data Security and Access}
\footnotesize
Suspicious access to KYC files outside regular business hours requires investigation.
\end{block}

\vspace{0.03cm}

\stepbox{\footnotesize
\steptitle{Step 1:} Understand the scenario. System logs show repeated access to KYC files outside regular business hours by an administrative account. \\
\steptitle{Step 2:} Recognize the red flags. \\
\quad - Outside hours: Unusual timing \\
\quad - Repeated access: Not a one-time event \\
\quad - Administrative account: Has privileged access \\
\steptitle{Step 3:} Understand the correct action. \textbf{Investigate the account's activity and suspend its privileges during the review}. \\
\steptitle{Step 4:} See why this is the right approach. \\
\quad - Investigation: Determines if data was compromised \\
\quad - Suspension: Prevents further unauthorized access \\
\steptitle{Step 5:} Apply the lesson. Suspicious access requires investigation and immediate suspension of privileges.}

\vspace{0.08cm}
\tipbox{\footnotesize \textbf{CORE LESSON:} Suspicious KYC access outside hours requires investigating the account and suspending privileges during the review.}
\end{frame}

% ================================================================
% SLIDE 35: TUTORIAL 34 — AI GOVERNANCE REGULATIONS
% ================================================================
\begin{frame}
\frametitle{\fontsize{14}{17}\selectfont TUTORIAL 34: AI Governance in Different Regions}

\small
\begin{block}{\fontsize{11}{13}\selectfont Topic: AI and AML Compliance}
\footnotesize
AI governance varies across regions. Understanding these differences is essential for global compliance.
\end{block}

\vspace{0.03cm}

\stepbox{\footnotesize
\steptitle{Step 1:} Understand that AI governance is not uniform. Different regions have different approaches. \\
\steptitle{Step 2:} Learn the regulatory features. Three key features are: \\
\quad (1) \textbf{Specific AI use-case restrictions} \\
\quad (2) \textbf{Requirements for human-in-the-loop oversight} \\
\quad (3) \textbf{Differing standards for transparency and data usage} \\
\steptitle{Step 3:} Understand why these differ across regions. \\
\quad - Different legal traditions \\
\quad - Different cultural values \\
\quad - Different risk appetites \\
\steptitle{Step 4:} Apply the lesson. Global institutions must understand and comply with regional AI regulations.}

\vspace{0.08cm}
\tipbox{\footnotesize \textbf{CORE LESSON:} AI governance varies by region with differences in use-case restrictions, human oversight requirements, and transparency standards.}
\end{frame}

% ================================================================
% SLIDE 36: TUTORIAL 35 — UAE ML ALGORITHMS
% ================================================================
\begin{frame}
\frametitle{\fontsize{14}{17}\selectfont TUTORIAL 35: Machine Learning in UAE AML}

\small
\begin{block}{\fontsize{11}{13}\selectfont Topic: UAE AML Technology}
\footnotesize
Machine learning algorithms provide specific advantages in the UAE's AML compliance context.
\end{block}

\vspace{0.03cm}

\stepbox{\footnotesize
\steptitle{Step 1:} Understand the UAE context. The UAE has a complex financial system with significant cross-border flows. \\
\steptitle{Step 2:} Learn the advantage of machine learning. \textbf{They improve customer risk assessments by identifying transaction anomalies}. \\
\steptitle{Step 3:} Understand why this is valuable. \\
\quad - Traditional systems miss complex patterns \\
\quad - Machine learning detects anomalies \\
\quad - Better risk assessment = better compliance \\
\steptitle{Step 4:} Apply the lesson. Machine learning improves AML compliance by identifying anomalies in customer transactions.}

\vspace{0.08cm}
\tipbox{\footnotesize \textbf{CORE LESSON:} Machine learning improves UAE AML compliance by identifying transaction anomalies for better customer risk assessments.}
\end{frame}

% ================================================================
% SLIDE 37: TUTORIAL 36 — ESG-ALIGNED COMPLIANCE
% ================================================================
\begin{frame}
\frametitle{\fontsize{14}{17}\selectfont TUTORIAL 36: ESG-Aligned Compliance Requirements}

\small
\begin{block}{\fontsize{11}{13}\selectfont Topic: ESG and Compliance}
\footnotesize
ESG-aligned compliance requirements are becoming increasingly important. Understanding these is essential.
\end{block}

\vspace{0.03cm}

\stepbox{\footnotesize
\steptitle{Step 1:} Understand what ESG-aligned compliance means. It integrates Environmental, Social, and Governance factors into compliance. \\
\steptitle{Step 2:} Learn the reporting efforts. Three examples are: \\
\quad (1) \textbf{Publishing carbon emissions by facility} \\
\quad (2) \textbf{Reporting executive salary data across genders} \\
\quad (3) \textbf{Reporting supply chain labor standards} \\
\steptitle{Step 3:} Understand why these are ESG-aligned. \\
\quad - Carbon emissions: Environmental \\
\quad - Salary data: Social \\
\quad - Labor standards: Social/Governance \\
\steptitle{Step 4:} Apply the lesson. ESG-aligned compliance includes environmental, social, and governance reporting.}

\vspace{0.08cm}
\tipbox{\footnotesize \textbf{CORE LESSON:} ESG-aligned compliance includes reporting on carbon emissions, executive salary data, and supply chain labor standards.}
\end{frame}

% ================================================================
% SLIDE 38: TUTORIAL 37 — ONLINE GAMBLING PARTNER DUE DILIGENCE
% ================================================================
\begin{frame}
\frametitle{\fontsize{14}{17}\selectfont TUTORIAL 37: Online Gambling Partner Due Diligence}

\small
\begin{block}{\fontsize{11}{13}\selectfont Topic: Third-Party Risk Management}
\footnotesize
Before processing payments for online gambling platforms, enhanced due diligence is required.
\end{block}

\vspace{0.03cm}

\stepbox{\footnotesize
\steptitle{Step 1:} Understand the scenario. A UK-based financial institution wants to process payments for an online gambling platform headquartered in another EEA country. \\
\steptitle{Step 2:} Recognize the risks. Online gambling is: \\
\quad - High-risk for money laundering \\
\quad - Subject to different regulations across countries \\
\quad - Often targeted by criminals \\
\steptitle{Step 3:} Understand the correct action. \textbf{Conduct enhanced due diligence and evaluate the AML framework of the EEA-based partner}. \\
\steptitle{Step 4:} See why this is required. \\
\quad - Ensure the partner has AML controls \\
\quad - Understand the regulatory context \\
\quad - Assess the risk \\
\steptitle{Step 5:} Apply the lesson. Before processing gambling payments, conduct EDD and evaluate the partner's AML framework.}

\vspace{0.08cm}
\tipbox{\footnotesize \textbf{CORE LESSON:} Before processing payments for online gambling, conduct enhanced due diligence and evaluate the partner's AML framework.}
\end{frame}

% ================================================================
% SLIDE 39: TUTORIAL 38 — PRIVATE BANKING RISK FACTORS
% ================================================================
\begin{frame}
\frametitle{\fontsize{14}{17}\selectfont TUTORIAL 38: Private Banking Risk Factors}

\small
\begin{block}{\fontsize{11}{13}\selectfont Topic: Private Banking Risks}
\footnotesize
Private banking has specific risk factors that increase money laundering risk.
\end{block}

\vspace{0.03cm}

\stepbox{\footnotesize
\steptitle{Step 1:} Understand the private banking context. It involves complex services for high-net-worth clients. \\
\steptitle{Step 2:} Recognize the risk factor. \textbf{Offshore jurisdictions used to establish client accounts} increases money laundering risk. \\
\steptitle{Step 3:} Understand why this creates risk. \\
\quad - Offshore jurisdictions: Have weak controls \\
\quad - Establish accounts: Creates distance from oversight \\
\quad - Harder to trace: More complex structures \\
\steptitle{Step 4:} Apply the lesson. Offshore jurisdiction accounts in private banking are a significant risk factor.}

\vspace{0.08cm}
\tipbox{\footnotesize \textbf{CORE LESSON:} Offshore jurisdictions used to establish client accounts are a key money laundering risk factor in private banking.}
\end{frame}

% ================================================================
% SLIDE 40: TUTORIAL 39 — INCOMPLETE WIRE FIELDS
% ================================================================
\begin{frame}
\frametitle{\fontsize{14}{17}\selectfont TUTORIAL 39: Incomplete Wire Transfer Fields}

\small
\begin{block}{\fontsize{11}{13}\selectfont Topic: Wire Transfer Compliance}
\footnotesize
Wires with incomplete name fields and abbreviated bank information should be returned.
\end{block}

\vspace{0.03cm}

\stepbox{\footnotesize
\steptitle{Step 1:} Understand the scenario. A series of wire transfers from a correspondent bank have incomplete name fields and abbreviated originating bank information. \\
\steptitle{Step 2:} Recognize the problem. Incomplete information violates wire transfer protocols and hinders traceability. \\
\steptitle{Step 3:} Understand the correct action. \textbf{Return the wire to notify the correspondent bank of formatting issues that violate wire transfer protocols}. \\
\steptitle{Step 4:} See why this is the right approach. \\
\quad - Returns enforce compliance \\
\quad - Notifies the bank of the issue \\
\quad - Protects against AML risk \\
\steptitle{Step 5:} Apply the lesson. Return incomplete wires to enforce protocol compliance.}

\vspace{0.08cm}
\tipbox{\footnotesize \textbf{CORE LESSON:} Return wires with incomplete fields to enforce wire transfer protocols and maintain traceability.}
\end{frame}

% ================================================================
% SLIDE 41: TUTORIAL 40 — BASEL AML INDEX
% ================================================================
\begin{frame}
\frametitle{\fontsize{14}{17}\selectfont TUTORIAL 40: Basel AML Index}

\small
\begin{block}{\fontsize{11}{13}\selectfont Topic: AML Risk Assessment Tools}
\footnotesize
The Basel AML Index provides valuable insights into money laundering vulnerabilities.
\end{block}

\vspace{0.03cm}

\stepbox{\footnotesize
\steptitle{Step 1:} Understand what the Basel AML Index is. It is a comprehensive index that measures country-level money laundering risk. \\
\steptitle{Step 2:} Recognize why it is valuable. \textbf{It compiles vulnerability indicators across legal, political, and financial dimensions}. \\
\steptitle{Step 3:} Understand what it covers. \\
\quad - Legal: Laws and enforcement \\
\quad - Political: Corruption and governance \\
\quad - Financial: Financial system vulnerabilities \\
\steptitle{Step 4:} Apply the lesson. The Basel AML Index provides multi-dimensional risk assessment for policy reform.}

\vspace{0.08cm}
\tipbox{\footnotesize \textbf{CORE LESSON:} The Basel AML Index compiles vulnerability indicators across legal, political, and financial dimensions for risk assessment.}
\end{frame}

% ================================================================
% SLIDE 42: TUTORIAL 41 — FATF IMMEDIATE OUTCOME 5
% ================================================================
\begin{frame}
\frametitle{\fontsize{14}{17}\selectfont TUTORIAL 41: FATF Immediate Outcome 5}

\small
\begin{block}{\fontsize{11}{13}\selectfont Topic: FATF Immediate Outcomes}
\footnotesize
FATF Immediate Outcome 5 addresses the misuse of legal persons for money laundering.
\end{block}

\vspace{0.03cm}

\stepbox{\footnotesize
\steptitle{Step 1:} Understand the scenario. A country has a high volume of legal arrangements relying on anonymous shell entities. Authorities struggle to trace beneficial ownership. \\
\steptitle{Step 2:} Recognize what this indicates. This situation indicates a weak rating for \textbf{Immediate Outcome 5 – Legal persons and arrangements are not misused}. \\
\steptitle{Step 3:} Understand what Immediate Outcome 5 covers. It assesses whether legal persons and arrangements are being misused for money laundering. \\
\steptitle{Step 4:} See why this is weak. \\
\quad - Anonymous shell entities: Hide ownership \\
\quad - Cannot trace beneficial ownership: Control is weak \\
\steptitle{Step 5:} Apply the lesson. Immediate Outcome 5 focuses on preventing misuse of legal persons.}

\vspace{0.08cm}
\tipbox{\footnotesize \textbf{CORE LESSON:} FATF Immediate Outcome 5 assesses whether legal persons and arrangements are misused. Anonymous shell entities indicate weakness.}
\end{frame}

% ================================================================
% SLIDE 43: TUTORIAL 42 — G-20 ACTION PLANS
% ================================================================
\begin{frame}
\frametitle{\fontsize{14}{17}\selectfont TUTORIAL 42: G-20 Anti-Corruption Action Plans}

\small
\begin{block}{\fontsize{11}{13}\selectfont Topic: G-20 Anti-Corruption}
\footnotesize
The G-20 Anti-Corruption Working Group coordinates international anti-corruption efforts.
\end{block}

\vspace{0.03cm}

\stepbox{\footnotesize
\steptitle{Step 1:} Understand the G-20 ACWG's role. It coordinates anti-corruption efforts among member states. \\
\steptitle{Step 2:} Learn how it contributes. It does this by \textbf{coordinating with member states to establish and implement multi-year anti-corruption action plans}. \\
\steptitle{Step 3:} Understand why this is effective. \\
\quad - Multi-year plans: Provide sustained effort \\
\quad - Coordination: Ensures consistency \\
\quad - Action plans: Provide clear objectives \\
\steptitle{Step 4:} Apply the lesson. The G-20 ACWG uses multi-year action plans to coordinate anti-corruption efforts.}

\vspace{0.08cm}
\tipbox{\footnotesize \textbf{CORE LESSON:} The G-20 ACWG coordinates with member states to establish and implement multi-year anti-corruption action plans.}
\end{frame}

% ================================================================
% SLIDE 44: TUTORIAL 43 — BENEFICIAL OWNERSHIP VERIFICATION
% ================================================================
\begin{frame}
\frametitle{\fontsize{14}{17}\selectfont TUTORIAL 43: Beneficial Ownership Verification}

\small
\begin{block}{\fontsize{11}{13}\selectfont Topic: FATF Recommendations 24 and 25}
\footnotesize
Technical compliance with beneficial ownership recommendations is not enough. Effectiveness requires verification mechanisms.
\end{block}

\vspace{0.03cm}

\stepbox{\footnotesize
\steptitle{Step 1:} Understand the scenario. A country receives technical compliance with Recommendations 24 and 25 but scores poorly on Immediate Outcome 5. \\
\steptitle{Step 2:} Recognize the gap. Technical compliance (having the laws) is not the same as effectiveness (making them work). \\
\steptitle{Step 3:} Understand the strategy to improve. \textbf{Develop verification mechanisms for beneficial ownership data, enforce disclosure requirements, and increase cooperation with foreign counterparts}. \\
\steptitle{Step 4:} See why this works. \\
\quad - Verification: Ensures data is accurate \\
\quad - Enforcement: Ensures compliance \\
\quad - Cooperation: Enables cross-border action \\
\steptitle{Step 5:} Apply the lesson. Technical compliance must be supported by verification, enforcement, and cooperation.}

\vspace{0.08cm}
\tipbox{\footnotesize \textbf{CORE LESSON:} Improve beneficial ownership effectiveness with verification mechanisms, enforcement of disclosure, and international cooperation.}
\end{frame}

% ================================================================
% SLIDE 45: TUTORIAL 44 — TERRORIST FINANCING COORDINATION
% ================================================================
\begin{frame}
\frametitle{\fontsize{14}{17}\selectfont TUTORIAL 44: Counter-Terrorist Financing Coordination}

\small
\begin{block}{\fontsize{11}{13}\selectfont Topic: Counter-Terrorist Financing}
\footnotesize
International information sharing is essential for effective counter-terrorist financing efforts.
\end{block}

\vspace{0.03cm}

\stepbox{\footnotesize
\steptitle{Step 1:} Understand the scenario. A cross-border team using scattered national legislation and ad hoc information channels failed to produce effective outcomes. \\
\steptitle{Step 2:} Recognize the key shortcoming. The \textbf{absence of coordinated international information sharing frameworks} was likely the key shortcoming. \\
\steptitle{Step 3:} Understand why this matters. \\
\quad - Terrorism is cross-border \\
\quad - Information must flow across borders \\
\quad - Ad hoc channels are insufficient \\
\steptitle{Step 4:} Apply the lesson. Effective CTF requires coordinated international information sharing, not ad hoc channels.}

\vspace{0.08cm}
\tipbox{\footnotesize \textbf{CORE LESSON:} Counter-terrorist financing requires coordinated international information sharing frameworks, not ad hoc channels.}
\end{frame}

% ================================================================
% SLIDE 46: TUTORIAL 45 — BULK CASH AND DECENTRALIZED STORAGE
% ================================================================
\begin{frame}
\frametitle{\fontsize{14}{17}\selectfont TUTORIAL 45: Bulk Cash and Decentralized Storage}

\small
\begin{block}{\fontsize{11}{13}\selectfont Topic: Terrorist Financing Vulnerabilities}
\footnotesize
Terrorists use bulk cash movement and decentralized storage to move money undetected.
\end{block}

\vspace{0.03cm}

\stepbox{\footnotesize
\steptitle{Step 1:} Understand the scenario. Attackers relied on an informal courier network to move cash across borders and stored money in small hotel safes. \\
\steptitle{Step 2:} Recognize what this highlights. This highlights \textbf{vulnerabilities in bulk cash movement and decentralized storage methods}. \\
\steptitle{Step 3:} Understand why these are vulnerabilities. \\
\quad - Bulk cash: No electronic trail \\
\quad - Decentralized storage: Hard to find \\
\quad - Courier networks: Informal, hard to track \\
\steptitle{Step 4:} Apply the lesson. Bulk cash and decentralized storage are significant vulnerabilities for terrorist financing.}

\vspace{0.08cm}
\tipbox{\footnotesize \textbf{CORE LESSON:} Bulk cash movement and decentralized storage methods are significant vulnerabilities that terrorists exploit.}
\end{frame}

% ================================================================
% SLIDE 47: TUTORIAL 46 — HUMAN TRAFFICKING RISKS
% ================================================================
\begin{frame}
\frametitle{\fontsize{14}{17}\selectfont TUTORIAL 46: Human Trafficking AML Risks}

\small
\begin{block}{\fontsize{11}{13}\selectfont Topic: Human Trafficking and AML}
\footnotesize
Human trafficking presents specific risks in the AML context. Understanding these is essential.
\end{block}

\vspace{0.03cm}

\stepbox{\footnotesize
\steptitle{Step 1:} Understand the risks. Two key risks are: \\
\quad (1) \textbf{Variability in how jurisdictions classify exploitation} \\
\quad (2) \textbf{Use of trafficked individuals for value transfer} \\
\steptitle{Step 2:} Understand why variability is a risk. Different jurisdictions classify exploitation differently. This creates gaps that traffickers exploit. \\
\steptitle{Step 3:} Understand how trafficked individuals are used. Traffickers use victims to move value — sometimes without their knowledge. \\
\steptitle{Step 4:} Apply the lesson. Human trafficking risks in AML include jurisdictional classification differences and use of victims for value transfer.}

\vspace{0.08cm}
\tipbox{\footnotesize \textbf{CORE LESSON:} Human trafficking AML risks include variability in jurisdiction classification of exploitation and use of victims for value transfer.}
\end{frame}

% ================================================================
% SLIDE 48: TUTORIAL 47 — OTC DERIVATIVES TRANSPARENCY
% ================================================================
\begin{frame}
\frametitle{\fontsize{14}{17}\selectfont TUTORIAL 47: OTC Derivatives Transparency}

\small
\begin{block}{\fontsize{11}{13}\selectfont Topic: OTC Derivatives Surveillance}
\footnotesize
Over-the-counter derivatives present transparency challenges. Specific actions improve transparency.
\end{block}

\vspace{0.03cm}

\stepbox{\footnotesize
\steptitle{Step 1:} Understand the scenario. A capital markets firm identifies gaps in post-trade surveillance for OTC derivatives. \\
\steptitle{Step 2:} Learn the actions that improve transparency. Two key actions are: \\
\quad (1) \textbf{Implement end-of-day reporting protocols} \\
\quad (2) \textbf{Enable automated reconciliation between front and back office systems} \\
\steptitle{Step 3:} Understand why these work. \\
\quad - End-of-day reporting: Provides daily transparency \\
\quad - Automated reconciliation: Catches discrepancies \\
\steptitle{Step 4:} Apply the lesson. Improve OTC derivative transparency through reporting protocols and automated reconciliation.}

\vspace{0.08cm}
\tipbox{\footnotesize \textbf{CORE LESSON:} Improve OTC derivatives transparency with end-of-day reporting protocols and automated reconciliation between front and back office.}
\end{frame}

% ================================================================
% SLIDE 49: SUMMARY
% ================================================================
\begin{frame}
\frametitle{\fontsize{14}{17}\selectfont SUMMARY: Key Lessons from All Topics}

\small
\begin{block}{\fontsize{11}{13}\selectfont What We Have Learned:}
\footnotesize
\begin{enumerate}
    \item \textbf{Open-Source Research:} Use to investigate unfamiliar entities.
    \item \textbf{Manual Escalation:} Frontline staff must escalate even without system alerts.
    \item \textbf{Documentation:} Use structured templates and adapt to jurisdictional requirements.
    \item \textbf{Bypassing Approvals:} Suggests potential layering activity.
    \item \textbf{E-Commerce Mitigation:} Use tiered onboarding and sanctions screening.
    \item \textbf{Token Risks:} Decentralized platforms obscure source and destination.
    \item \textbf{Recommendation 18:} Requires group-wide AML programs.
    \item \textbf{FSRB Role:} Implement and promote FATF standards regionally.
    \item \textbf{IOSCO Objectives:} Investor protection and financial stability.
    \item \textbf{IMF Evaluations:} Economic improvements can mask financial crime.
    \item \textbf{False Positives:} Tune systems, don't reduce criteria.
    \item \textbf{Adverse Media:} Automation improves efficiency.
    \item \textbf{AI Inconsistency:} Engage governance committee.
    \item \textbf{Device Intelligence:} Captures device-specific information.
    \item \textbf{Wire Investigations:} Investigate wires through small foreign banks.
    \item \textbf{TBML Controls:} Digitized bills, dual approvals, look-back periods.
    \item \textbf{Synthetic Fraud:} Suspend product, update risk assessment.
    \item \textbf{Private Banking Risks:} Numbered accounts, discretion, high assets.
    \item \textbf{Egmont Group:} Secure communication, shared typologies.
    \item \textbf{FATF Response:} Training and tighter oversight.
    \item \textbf{EDD for Common Owner:} Conduct enhanced due diligence.
    \item \textbf{Concentration Accounts:} Internal cash movements and back-office.
    \item \textbf{PEP Management:} Detection, verification, monitoring.
    \item \textbf{Prepaid Cards:} Anonymous, especially reloadable.
    \item \textbf{Sanctions Lists:} FATF, UN, State Department.
    \item \textbf{Data Sources:} Internal = behavior, External = lists.
    \item \textbf{Data Lineage:} Transformations and source capture.
    \item \textbf{Basel Institute:} ICAR + capacity building.
    \item \textbf{World Bank/IMF:} Reform AML/CFT frameworks.
    \item \textbf{G-20 ACWG:} Asset recovery efficiency.
    \item \textbf{PEP Escalation:} Protocols and training.
    \item \textbf{EU Directive:} Post-2008 crisis updates.
    \item \textbf{KYC Access:} Investigate and suspend privileges.
    \item \textbf{AI Governance:} Varies by region.
    \item \textbf{UAE ML:} Anomaly detection.
    \item \textbf{ESG Compliance:} Carbon, salary, labor reporting.
    \item \textbf{Gambling EDD:} Evaluate partner AML framework.
    \item \textbf{Private Banking:} Offshore jurisdictions increase risk.
    \item \textbf{Incomplete Wires:} Return to enforce protocols.
    \item \textbf{Basel AML Index:} Legal, political, financial indicators.
    \item \textbf{Immediate Outcome 5:} Legal persons misuse.
    \item \textbf{G-20 Action Plans:} Multi-year coordination.
    \item \textbf{BO Verification:} Verification, enforcement, cooperation.
    \item \textbf{CTF Coordination:} International information sharing.
    \item \textbf{Bulk Cash:} Movement and decentralized storage vulnerabilities.
    \item \textbf{Human Trafficking:} Classification variability, value transfer.
    \item \textbf{OTC Transparency:} Reporting protocols, automated reconciliation.
\end{enumerate}
\end{block}

\vspace{0.05cm}
\tipbox{\footnotesize \textbf{FINAL LESSON:} Every topic teaches us that understanding the \textit{underlying principle} is more important than memorizing facts. Know \textit{why} and you will know \textit{what} to do.}
\end{frame}

% ================================================================
% SLIDE 50: THANK YOU
% ================================================================
\begin{frame}
\centering
\vspace{2.5cm}
{\fontsize{22}{26}\selectfont \textbf{Thank You}} \\[0.3cm]
{\fontsize{14}{17}\selectfont Keep learning. Keep growing.} \\[0.2cm]
\rule{3cm}{1pt} \\[0.2cm]
{\fontsize{10}{12}\selectfont \textit{Compliance is not about perfection — it is about continuous improvement.}}
\vspace{0.5cm}
{\fontsize{9}{11}\selectfont \textcolor{darkgray}{End of Tutorial}}
\end{frame}

\end{document}